\section{Detection Playbooks}

Detection playbooks translate knowledge of adversary behavior into concrete, repeatable logic that can run continuously in Splunk or Microsoft Sentinel. Unlike hunting playbooks, which are exploratory and analyst-driven, detection playbooks are designed for automation: they define what to look for, how to detect it, how to investigate it, and how to tune it over time. For a CTI analyst or detection engineer, these playbooks form the backbone of a mature detection program and serve as a long-term reference for building and maintaining high-fidelity alerts.

Each detection playbook in this section follows a consistent structure. It begins with a short description of the behavior and why it matters, followed by the required telemetry. Then it provides SPL and KQL detection logic, along with Linux investigation commands where relevant. Finally, it discusses tuning considerations and common sources of false positives. The goal is not only to provide queries, but to explain how and why they work, so that they can be adapted to different environments and threat landscapes.

\subsection*{Detection Playbook: Encoded PowerShell Commands}

\textbf{Behavior and Rationale}

Adversaries frequently abuse PowerShell to execute malicious code, download payloads, or perform reconnaissance. To evade simple string-based detections, they often use the \texttt{-EncodedCommand} parameter, which allows them to supply Base64-encoded commands. This behavior is highly suspicious in most environments and is a strong candidate for a high-priority detection, especially on servers and high-value endpoints.

\textbf{Required Telemetry}

This detection relies on detailed process creation logs that include command-line arguments. On Windows, Sysmon Event ID 1 or Windows Event ID 4688 with full command-line logging is ideal. In Splunk, these events are typically found in indexes such as \texttt{sysmon} or \texttt{wineventlog}. In Sentinel, the equivalent data is often in \texttt{Sysmon} or \texttt{SecurityEvent} tables.

\textbf{SPL Detection Logic}
\begin{lstlisting}
index=sysmon EventCode=1 Image="*powershell.exe"
| search CommandLine="*EncodedCommand*"
| stats count by Computer, User, ParentImage, Image, CommandLine
| where count >= 1
\end{lstlisting}

\textbf{KQL Detection Logic}
\begin{lstlisting}
Sysmon
| where EventID == 1 and Image endswith "powershell.exe"
| where CommandLine contains "EncodedCommand"
| summarize Count=count() by Computer, User, ParentImage, Image, CommandLine
| where Count >= 1
\end{lstlisting}

This logic looks for any PowerShell process that includes the \texttt{EncodedCommand} parameter. In many environments, even a single occurrence is worth investigating. The aggregation step (\texttt{stats} / \texttt{summarize}) helps group similar events and makes it easier to review repeated activity.

\textbf{Linux Investigation Pivot}

On systems with Sysmon for Linux or similar logging, analysts can pivot from SIEM detections to local logs:

\begin{lstlisting}[language=bash]
grep -i "powershell" /var/log/sysmon/sysmon.log | grep -i "EncodedCommand"
\end{lstlisting}

\textbf{Tuning and False Positives}

Some administrative scripts or management tools may legitimately use encoded PowerShell, especially in automation-heavy environments. To reduce noise, detections can be tuned by excluding known management accounts, specific hosts (such as configuration servers), or known-good scripts. Over time, a whitelist of expected encoded PowerShell usage can be maintained, while all other occurrences remain high-priority alerts.

\subsection*{Detection Playbook: Excessive Failed Logons}

\textbf{Behavior and Rationale}

Repeated failed logons are a classic indicator of brute-force attacks, password spraying, or credential guessing. While individual failures are common and often benign, patterns of repeated failures against specific accounts or from specific hosts are highly suspicious. Detecting these patterns early helps prevent account compromise and lateral movement.

\textbf{Required Telemetry}

On Windows, failed logons are recorded as Event ID 4625 in the Security log. In Splunk, these events are typically found in \texttt{wineventlog} or similar indexes. In Sentinel, they appear in the \texttt{SecurityEvent} table. For Linux, failed SSH logons are recorded in \texttt{/var/log/auth.log}.

\textbf{SPL Detection Logic}
\begin{lstlisting}
index=wineventlog EventCode=4625
| stats count by AccountName, WorkstationName
| where count > 10
\end{lstlisting}

\textbf{KQL Detection Logic}
\begin{lstlisting}
SecurityEvent
| where EventID == 4625
| summarize Count=count() by Account, Workstation
| where Count > 10
\end{lstlisting}

This detection identifies accounts that have more than ten failed logons from a given workstation. The threshold can be adjusted based on the environment’s normal behavior. In high-noise environments, it may be useful to add a time window (for example, failures within one hour) to focus on concentrated attack patterns.

\textbf{Linux Investigation Pivot}

On Linux systems, failed SSH logons can be reviewed directly:

\begin{lstlisting}[language=bash]
grep "Failed password" /var/log/auth.log | awk '{print $9}' | sort | uniq -c
\end{lstlisting}

This command shows how many times each user has experienced failed logons, which can be correlated with SIEM detections.

\textbf{Tuning and False Positives}

False positives often arise from misconfigured services, automated scripts with outdated credentials, or users repeatedly entering incorrect passwords. Tuning should focus on excluding known noisy systems and service accounts, while maintaining strict thresholds for privileged accounts and high-value targets. Additional context, such as geolocation or device type, can further improve fidelity.

\subsection*{Detection Playbook: Rare Parent-Child Process Relationships}

\textbf{Behavior and Rationale}

Legitimate processes tend to follow predictable parent-child relationships. For example, \texttt{explorer.exe} commonly spawns user applications, while \texttt{services.exe} spawns services. When an unusual parent process launches a sensitive child process (such as a scripting engine or system tool), it may indicate process injection, masquerading, or other malicious activity. Detecting rare parent-child combinations is a powerful behavioral technique.

\textbf{Required Telemetry}

This detection relies on detailed process creation logs with parent and child process information. Sysmon Event ID 1 is ideal, as it includes \texttt{ParentImage} and \texttt{Image} fields. In Splunk, these events are typically stored in a \texttt{sysmon} index; in Sentinel, they appear in the \texttt{Sysmon} table.

\textbf{SPL Detection Logic}
\begin{lstlisting}
index=sysmon EventCode=1
| stats count by ParentImage, Image
| where count < 3
\end{lstlisting}

\textbf{KQL Detection Logic}
\begin{lstlisting}
Sysmon
| where EventID == 1
| summarize Count=count() by ParentImage, Image
| where Count < 3
\end{lstlisting}

This detection highlights parent-child pairs that occur fewer than three times in the dataset. These rare combinations are strong candidates for further investigation. In practice, this detection is often run over a rolling time window (for example, the last seven days) to reflect current behavior.

\textbf{Linux Investigation Pivot}

On Linux, process relationships can be inspected using \texttt{ps} and \texttt{pstree}:

\begin{lstlisting}[language=bash]
ps aux --forest
\end{lstlisting}

This view helps analysts confirm whether a suspicious process tree observed in SIEM logs is legitimate or malicious.

\textbf{Tuning and False Positives}

Rare does not always mean malicious. New software deployments, updates, or one-time administrative tasks can create unusual process chains. Tuning involves building a baseline of known-good rare pairs and excluding them from future alerts. Over time, the detection should focus on truly unexpected combinations, especially those involving scripting engines, interpreters, or administrative tools.

\subsection*{Detection Playbook: Rare Outbound Network Connections}

\textbf{Behavior and Rationale}

Compromised hosts often communicate with external command-and-control servers or exfiltration endpoints that are not part of normal business operations. These connections are typically low-frequency and may target unusual ports or IP addresses. Detecting rare outbound connections helps identify such activity, especially when combined with CTI enrichment.

\textbf{Required Telemetry}

This detection relies on network connection logs that include destination IPs and ports. Sysmon Event ID 3 provides this information on Windows. In Splunk, these events are often stored in a \texttt{sysmon} or network index; in Sentinel, they appear in the \texttt{Sysmon} table or dedicated network tables.

\textbf{SPL Detection Logic}
\begin{lstlisting}
index=sysmon EventCode=3
| stats count by DestinationIp, DestinationPort
| where count < 5
\end{lstlisting}

\textbf{KQL Detection Logic}
\begin{lstlisting}
Sysmon
| where EventID == 3
| summarize Count=count() by DestinationIp, DestinationPort
| where Count < 5
\end{lstlisting}

This detection surfaces destination IP and port combinations that appear fewer than five times in the dataset. These rare connections can then be enriched with CTI (for example, checking against threat feeds) to determine whether they are associated with known malicious infrastructure.

\textbf{Linux Investigation Pivot}

On Linux, active connections can be inspected with:

\begin{lstlisting}[language=bash]
ss -tulnp
\end{lstlisting}

This command shows listening and established connections, along with associated processes, which can be correlated with SIEM detections.

\textbf{Tuning and False Positives}

False positives may include legitimate but infrequent connections, such as one-off vendor support sessions or new cloud services. Tuning should focus on excluding known business destinations and internal infrastructure, while maintaining strict scrutiny for unknown external IPs, unusual ports, and connections from high-value systems.

\subsection*{Detection Playbook: Persistence via Scheduled Tasks and Cron Jobs}

\textbf{Behavior and Rationale}

Persistence mechanisms allow adversaries to maintain access to a system even after reboots or user logouts. On Windows, scheduled tasks are a common persistence technique; on Linux, cron jobs serve a similar purpose. Detecting the creation or modification of these mechanisms is critical for identifying long-term footholds.

\textbf{Required Telemetry}

On Windows, scheduled task creation is recorded as Event ID 4698 in the Security log and may also appear in Task Scheduler logs. In Splunk, these events are typically found in \texttt{wineventlog}. In Sentinel, they appear in the \texttt{SecurityEvent} table. On Linux, cron configurations are stored in \texttt{/etc/crontab}, \texttt{/etc/cron.d/}, and user-specific cron files.

\textbf{SPL Detection Logic}
\begin{lstlisting}
index=wineventlog EventCode=4698
| table _time, Computer, TaskName, Author, Command
\end{lstlisting}

\textbf{KQL Detection Logic}
\begin{lstlisting}
SecurityEvent
| where EventID == 4698
| project TimeGenerated, Computer, TaskName, SubjectUserName, Command
\end{lstlisting}

This detection surfaces newly created scheduled tasks, including the task name, author, and command. In many environments, scheduled task creation is relatively rare and should be closely monitored, especially on servers and high-value endpoints.

\textbf{Linux Investigation Pivot}

On Linux, cron-based persistence can be reviewed with:

\begin{lstlisting}[language=bash]
cat /etc/crontab
ls -al /etc/cron.d/
crontab -l
\end{lstlisting}

These commands reveal system-wide and user-specific cron jobs, which can be compared against known-good configurations.

\textbf{Tuning and False Positives}

Legitimate administrative tasks, backup jobs, and monitoring scripts often use scheduled tasks or cron jobs. Tuning should focus on identifying expected tasks (by name, author, or command) and excluding them from alerts. New or modified tasks that execute scripting engines, interpreters, or binaries from unusual paths should remain high-priority.

\subsection*{Summary}

The detection playbooks in this section illustrate how to move from abstract knowledge of adversary behavior to concrete, automated logic in Splunk and Sentinel. Each playbook combines behavioral focus, clear telemetry requirements, SPL and KQL implementations, Linux investigation pivots, and practical tuning guidance. Over time, these playbooks can be expanded, refined, and adapted to new threats, forming a living detection library that supports your long-term growth as a CTI analyst, threat hunter, and detection engineer.
